\documentclass{article}
\usepackage{amsmath}
\usepackage{amssymb}
\usepackage{amsthm}
\usepackage{enumitem}
\usepackage[a4paper, margin=1in]{geometry}

\newtheorem{theorem}{Teorema}
\newtheorem{corollary}{Corollario}

\begin{document}
\setlength{\parindent}{0pt}

\section{Matrici}
\subsection{Definizione di matrice}
    Una matrice $A$ di dimensioni $m \times n$ a coefficienti reali è un oggetto matematico della forma:
    \[  A =
        \begin{bmatrix}
            a_{11} & a_{12} & ... & a_{1n} \\
            a_{21} & a_{22} & ... & a_{2n} \\
            \vdots & \ddots & \ddots & \vdots \\
            a_{m1} & a_{m2} & ... & a_{mn} 
        \end{bmatrix}
    \]
    Dove $a_{ij} \in \mathbb{R}$ sono i coefficienti della matrice, $m$ è il numero delle righe ed $n$ il numero delle colonne.
    Se $n = m$ la matrice si dice quadrata e gli elementi $\{ a_{ii} \; | \; i \in \{0, ..., n\}\}$ sono chiamati diagonale della matrice. 

    \vspace{10pt}
    Sia $\mathbb{R}^{m, n}$ l'insieme di tutte le matrici $m \times n$, possiamo affermare che sussiste una biezione fra $\mathbb{R}^{m, 1} \leftrightarrow \mathbb{R}^{m}$ e fra $\mathbb{R}^{1, n} \leftrightarrow \mathbb{R}^{n}$.



\subsection{Somma fra matrici}
    Siano $A, B \in \mathbb{R}^{m, n}$ due matrici $m \times n$, la somma di $A$ con $B$ è definita come:
    \[
        A + B = 
        \begin{bmatrix}
            a_{11} & a_{12} & ... & a_{1n} \\
            a_{21} & a_{22} & ... & a_{2n} \\
            \vdots & \ddots & \ddots & \vdots \\
            a_{m1} & a_{m2} & ... & a_{mn} 
        \end{bmatrix}
        +
        \begin{bmatrix}
            b_{11} & b_{12} & ... & b_{1n} \\
            b_{21} & b_{22} & ... & b_{2n} \\
            \vdots & \ddots & \ddots & \vdots \\
            b_{m1} & b_{m2} & ... & b_{mn} 
        \end{bmatrix}
        =
        \begin{bmatrix}
            a_{11} + b_{11} & a_{12} + b_{12} & ... & a_{1n} + b_{1n} \\
            a_{21} + b_{21} & a_{22} + b_{22} & ... & a_{2n} + b_{2n} \\
            \vdots & \ddots & \ddots & \vdots \\
            a_{m1} + b_{m1} & a_{m2} + b_{m2} & ... & a_{mn} + b_{mn} 
        \end{bmatrix}
    \]

    \vspace{10pt}
    Proprietà:
    \begin{enumerate}[label=(\arabic*)]
        \item Commutatività: $\forall A, B \in \mathbb{R}^{m, n}, \; A + B = B + A$
        \item Associatività: $\forall A, B, C \in \mathbb{R}^{m, n}, \; (A + B) + C = A + (C + B)$
        \item Elemento neutro: $\exists N \in \mathbb{R}^{m, n} \; | \; \forall A \in \mathbb{R}^{m, n}, \; A + N = A $ 
        \item Elemento opposto: $\forall A \in \mathbb{R}^{m, n}, \; \exists -A \in \mathbb{R}^{m, n} \; | \; A + (-A) = N$
    \end{enumerate}



\subsection{Prodotto scalare-matrice}
    Sia $A \in \mathbb{R}^{m, n}$ una matrice e sia $\lambda \in \mathbb{R}$ uno scalare, il prodotto scalare-matrice di $\lambda$ con $A$ è definito come:
    \[
    \lambda A = 
    \begin{bmatrix}
        \lambda a_{11} & \lambda a_{12} & ... & \lambda a_{1n} \\
        \lambda a_{21} & \lambda a_{22} & ... & \lambda a_{2n} \\
        \vdots & \ddots & \ddots & \vdots \\
        \lambda a_{m1} & \lambda a_{m2} & ... & \lambda a_{mn} 
    \end{bmatrix} 
    \]

    \vspace{10pt}
    Proprietà:
    \begin{enumerate}[label=(\arabic*)]
        \item Associatività: $\forall \lambda, \mu \in \mathbb{R}, \; \forall A \in \mathbb{R}^{m, n}, \; (\lambda \mu) A = \lambda (\mu A)$ 
        \item Elemento neutro: $\forall A \in \mathbb{R}^{m, n} \; \exists 1 \in \mathbb{R} \; | \; 1A = A$
        \item Distributività rispetto alla somma fra matrici: $\forall \lambda \in \mathbb{R}, \; \forall A, B \in \mathbb{R}^{m, n}, \; \lambda(A + B) = \lambda A + \lambda B$
        \item Distributività rispetto alla somma fra scalari: $\forall \lambda, \mu \in \mathbb{R}, \; \forall A \in \mathbb{R}^{m, n}, \; (\lambda + \mu)A = \lambda A + \mu A$
    \end{enumerate}



\subsection{Matrici trasposte}
    Data una matrice $A \in \mathbb{R}^{m, n}$, la sua trasposta $A^{T} \in \mathbb{R}^{n, m}$ è la matrice nella forma:
    \[
        A = 
        \begin{bmatrix}
            a_{11} & a_{12} & ... & a_{1n} \\
            a_{21} & a_{22} & ... & a_{2n} \\
            \vdots & \ddots & \ddots & \vdots \\
            a_{m1} & a_{m2} & ... & a_{mn} 
        \end{bmatrix} 
        \implies
        A^{T} = 
        \begin{bmatrix}
            a_{11} & a_{21} & ... & a_{m1} \\
            a_{12} & a_{22} & ... & a_{m2} \\
            \vdots & \ddots & \ddots & \vdots \\
            a_{n1} & a_{n2} & ... & a_{nm} 
        \end{bmatrix} 
    \]


    Una matrice $A \in \mathbb{R}^{n, n}$ si dice simmetrica se $A = A^{T}$ oppure asimmetrica se $A = -A^{T}$

    \vspace{10pt}
    \begin{theorem}
        Sia $V$ lo spazio vettoriale determinato da $(\mathbb{R}^{n, n}, \mathbb{R}, +, \cdot)$, allora $S = \{A \in \mathbb{R}^{n, n} \; | \; A = A^{T}\}$ e $W = \{A \in \mathbb{R}^{n, n} \; | \; A = -A^{T}\}$ sono 
        sottospazi di $V$ e $S \bigoplus W = \mathbb{R}^{n, n}$ 
    \end{theorem}
    \begin{proof}[Dimostrazione]
        \begin{enumerate}[label=(\arabic*)]
            \item Dato che vale $\forall A_1, A_2 \in S, \; (A_1 + A_2)^{T} = A_1^T + A_2^T = A_1 + A_2$ e vale $\forall A \in S, \; (\lambda A)^T = \lambda A^T = \lambda A$ allora $S$ è sottospazio di $V$
            \item Dato che vale $\forall A_1, A_2 \in W, \; (A_1 + A_2)^{T} = A_1^T + A_2^T = -A_1 + -A_2 = -(A_1 + A_2)$ e vale $\forall A \in S, \; (\lambda A)^T = \lambda A^T = - \lambda A$ allora $W$ è sottospazio di $V$
            \item Se $A \in S \cap W$ allora $A = A^T = -A \implies A = 0_{\mathbb{R}^{n, n}}$, quindi la somma $S \bigoplus W$ è diretta.
            \item Sia $M \in \mathbb{R}^{n, n}$. Allora se $A_1 = \frac{M + M^{T}}{2} \in S$ e $A_2 = \frac{M - M^{T}}{2} \in W$ abbiamo che $A_1 + A_2 = M \in S \bigoplus W$
        \end{enumerate}
    \end{proof}

\subsection{Prodotto matrice-matrice}
    Date due matrici $A \in \mathbb{R}^{m, n}$ e $B \in \mathbb{R}^{n, k}$ il prodotto matrice-matrice fra $A$ e $B$ è definito come:
    \[
        A \cdot B = 
        \begin{bmatrix}
            a_{11} & a_{12} & ... & a_{1n} \\
            a_{21} & a_{22} & ... & a_{2n} \\
            \vdots & \ddots & \ddots & \vdots \\
            a_{m1} & a_{m2} & ... & a_{mn} 
        \end{bmatrix}
        \cdot 
        \begin{bmatrix}
            b_{11} & b_{12} & ... & b_{1k} \\
            b_{21} & b_{22} & ... & b_{2k} \\
            \vdots & \ddots & \ddots & \vdots \\
            b_{n1} & b_{n2} & ... & b_{nk} 
        \end{bmatrix}
        = 
        \begin{bmatrix}
            \displaystyle\sum_{i = 0}^{n}{a_{1i} b_{i1}} & \displaystyle\sum_{i = 0}^{n}{a_{1i} b_{i2}} & ... & \displaystyle\sum_{i = 0}^{n}{a_{1i} b_{ik}} \\
            \displaystyle\sum_{i = 0}^{n}{a_{2i} b_{i1}} & \displaystyle\sum_{i = 0}^{n}{a_{2i} b_{i2}} & ... & \displaystyle\sum_{i = 0}^{n}{a_{2i} b_{ik}} \\
            \vdots & \ddots & \ddots & \vdots \\
            \displaystyle\sum_{i = 0}^{n}{a_{mi} b_{i1}} & \displaystyle\sum_{i = 0}^{n}{a_{mi} b_{i2}} & ... & \displaystyle\sum_{i = 0}^{n}{a_{mi} b_{ik}} 
        \end{bmatrix}
        \in \mathbb{R}^{m, k}
    \]

    \vspace{10pt}
    Proprietà:
    \begin{enumerate}[label=(\arabic*)]
        \item Associatività: $\forall A \in \mathbb{R}^{m, n}, \forall B \in \mathbb{R}^{n, k}, \forall C \in \mathbb{R}^{k, p}, \; (A \cdot B) \cdot C = A \cdot (B \cdot C)$ 
        \item Associatività rispetto al prodotto scalare-matrice: $\forall \lambda \in \mathbb{R}, \forall A \in \mathbb{R}^{m, n}, \forall B \in \mathbb{R}^{n, k}, \; \lambda (A \cdot B) = \lambda A \cdot B = A \cdot \lambda B$
        \item Distributività sinistra rispetto alla somma fra matrici: $\forall A, B \in \mathbb{R}^{m, n}, \forall C \in \mathbb{R}^{k, m}, \; C \cdot (A + B) = C \cdot A + C \cdot B$
        \item Distributività destra rispetto alla somma fra matrici: $\forall A, B \in \mathbb{R}^{m, n}, \forall C \in \mathbb{R}^{n, k}, \; (A + B) \cdot C = C \cdot A + C \cdot B$
        \item Distributività della trasposizione rispetto al prodotto matrice-matrice: $\forall A \in \mathbb{R}^{m, n}, \forall B \in \mathbb{R}^{n, k}, \; (A \cdot B)^{T} = A^{T} \cdot B^{T}$
    \end{enumerate}

\subsection{Matrici ridotte, pivot e rango di matrice}
    Data una matrice $A \in \mathbb{R}^{m, n}$ essa si dice ridotta per riga se e solo se il primo elemento non nullo
    di ogni riga, chiamato pivot, si trova a destra del pivot della riga superiore. $A$ si dice invece fortemente ridotta,
    oppure in forma echelon (reduced row echelon form), se è ridotta per righe, ogni pivot vale 1 e il pivot è l'unico elemento non nullo della colonna. 

    \vspace{10pt}
    Un esempio di matrice ridotta per righe è:
    \[
        A = 
        \begin{bmatrix}
            3 & 2 & 4 & 9 \\
            0 & 1 & 8 & 8 \\
            0 & 0 & 0 & 6 \\
            0 & 0 & 0 & 0 
        \end{bmatrix} 
    \]
    Mentre un esempio di matrice fortemente ridotta per righe: 
    \[
        A = 
        \begin{bmatrix}
            1 & 0 & 0 & 9 \\
            0 & 1 & 0 & 8 \\
            0 & 0 & 1 & 6 \\
            0 & 0 & 0 & 0 
        \end{bmatrix} 
    \]

    \vspace{10pt}
    Ogni matrice $A \in \mathbb{R}^{m, n}$ può essere portata in forma fortemente ridotta per righe attraverso un processo iterativo chiamato
    eliminazione di Gauss, il quale consiste nell'applicare una delle seguenti traformazioni alle righe laddove è necessario:
    \begin{enumerate}[label=(E\textsubscript{\arabic*}), leftmargin=1cm]
        \item $R_i \rightarrow R_i + \lambda R_j$ con $\lambda \in \mathbb{R}$ 
        \item $R_i \rightarrow \lambda R_i$ con $\lambda \in \mathbb{R} - \{0\}$ 
        \item $R_i \rightarrow R_j$ 
    \end{enumerate}
    Dove $R_i, R_j \in \mathbb{R}^{1, n}$ sono le righe i e j-esime della matrice $A$.

    \vspace{10pt}
    Un esempio di eliminazione di Gauss è:
    \[
        A = 
        \begin{bmatrix}
            3 & 1 & 4 & 9 \\
            3 & 2 & 8 & 8 \\
            0 & 3 & 0 & 6 \\
            0 & 0 & 0 & 7 
        \end{bmatrix} 
        \xrightarrow{R_2 - R_1 \rightarrow R_2 \; (E_1)}
        \begin{bmatrix}
            3 & 1 & 4 & 9 \\
            0 & 1 & 4 & -1 \\
            0 & 3 & 0 & 6 \\
            0 & 0 & 0 & 7 
        \end{bmatrix} 
        \xrightarrow{R_3 - 3 R_2 \rightarrow R_3 \; (E_1)}
        \begin{bmatrix}
            3 & 1 & 4 & 9 \\
            0 & 1 & 4 & -1 \\
            0 & 0 & -12 & 9 \\
            0 & 0 & 0 & 7 
        \end{bmatrix} 
    \]
    
    \vspace{10pt}
    Il rango di una matrice è definito come il numero di righe (o colonne) linearmente indipendenti.

    \vspace{10pt}
    \begin{theorem}[Teorema del rango]
        Sia $A \in \mathbb{R}^{m, n}$ matrice e sia $R_A = \textnormal{span}(\mathbf{R_1}, \mathbf{R_2}, ..., \mathbf{R_m})$ lo spazio generato dalle righe di $A$, mentre $C_A = \textnormal{span}(\mathbf{C_1}, \mathbf{C_2}, ..., \mathbf{C_n})$ lo spazio
        generato dalle colonne di $A$, allora $\dim(R_A) = \dim(C_A)$     
    \end{theorem}

    \begin{proof}[Dimostrazione]
        Sia $B_R = \{\mathbf{r_1}, ..., \mathbf{r_k}\}$ una base di $R_A$, allora ogni riga $\mathbf{R_i} = a_{1}^{i} \mathbf{r_1} + a_{2}^{i} \mathbf{r_2} + ... + a_{k}^{i} \mathbf{r_k}$. Dunque ogni colonna 
        \[
            \mathbf{C_i} = 
            \begin{bmatrix}
                a_{1}^{1} \\
                a_{1}^{2} \\
                \vdots \\
                a_{1}^{m}
            \end{bmatrix}
            r_{1}^{i}
            + 
            \begin{bmatrix}
                a_{2}^{1} \\
                a_{2}^{2} \\
                \vdots \\
                a_{2}^{m}
            \end{bmatrix}
            r_{2}^{i}
            +
            \dots
            +
            \begin{bmatrix}
                a_{k}^{1} \\
                a_{k}^{2} \\
                \vdots \\
                a_{k}^{m}
            \end{bmatrix}
            r_{k}^{i}
        \]
        
        il che implica che 
        $G_C = \{
            \begin{bmatrix}
                a_{1}^{1} \\
                a_{1}^{2} \\
                \vdots \\
                a_{1}^{m}
            \end{bmatrix},  
            \begin{bmatrix}
                a_{2}^{1} \\
                a_{2}^{2} \\
                \vdots \\
                a_{2}^{m}
            \end{bmatrix},
            ...,
            \begin{bmatrix}
                a_{k}^{1} \\
                a_{k}^{2} \\
                \vdots \\
                a_{k}^{m}
            \end{bmatrix} 
            \}$
        sia un generatore di $C_A$. Quindi, per il lemma di 

        \vspace{5pt}
        Steinitz, si ha che $\dim(C_A) \leq |G_C| = |B_R| = \dim(R_A)$. Si prenda in considerazione $A^T$ e si ripeta il procedimento svolto.
        Si giungerà alla conclusione che $\dim(C_{A^T}) = \dim(R_A) \leq \dim(R_{A^T}) = \dim(C_A)$, per cui $\dim(R_A) = \dim(C_A)$.
    \end{proof}

\vspace{10pt}
\subsection{Sistemi lineari}
    Un sistema lineare è un sistema di $m$ equazioni lineari a $n$ incognite della forma:
    \[
        \begin{cases}
            a_{11}x_1 + a_{12}x_2 + ... + a_{1n}x_n = b_1 \\
            a_{21}x_1 + a_{22}x_2 + ... + a_{2n}x_n = b_2 \\
            \vdots \\
            a_{m1}x_1 + a_{m2}x_2 + ... + a_{mn}x_n = b_m
        \end{cases} 
    \]
    Per ogni sistema lineare è possibile costruire una sua versione "in matrice", ovvero:
    \[
        \begin{cases}
            a_{11}x_1 + a_{12}x_2 + ... + a_{1n}x_n = b_1 \\
            a_{21}x_1 + a_{22}x_2 + ... + a_{2n}x_n = b_2 \\
            \vdots \\
            a_{m1}x_1 + a_{m2}x_2 + ... + a_{mn}x_n = b_m
        \end{cases} 
        \, 
        \rightarrow 
        \,
        \begin{bmatrix}
            a_{11} & a_{12} & ... & a_{1n} \\
            a_{21} & a_{22} & ... & a_{2n} \\
            \vdots & \ddots & \ddots & \vdots \\
            a_{m1} & a_{m2} & ... & a_{mn} 
        \end{bmatrix}
        \,
        \begin{bmatrix}
            x_{1} \\ 
            x_{2} \\
            \vdots \\
            x_{n} \\
        \end{bmatrix}
        \,
        = 
        \begin{bmatrix}
            b_{1} \\ 
            b_{2} \\
            \vdots \\
            b_{m} \\
        \end{bmatrix}
        \,
        \rightarrow
        A\mathbf{x} = \mathbf{b}
    \]
    La matrice aumentata del sistema $(A | B) \in \mathbb{R}^{m, n+1}$ è definita come:
    \[
        \left[
        \begin{array}{cccc | c}
            a_{11} & a_{12} & ... & a_{1n} & b_1 \\
            a_{21} & a_{22} & ... & a_{2n} & b_2 \\
            \vdots & \ddots & \ddots & \vdots & \vdots \\
            a_{m1} & a_{m2} & ... & a_{mn} & b_m 
        \end{array}
        \right]
    \]

    Un sistema lineare cui $\mathbf{b} = \mathbf{0}$ viene detto omogeneo ed è sempre risolubile in quanto ammette la soluzione triviale $\mathbf{x} = \mathbf{0}$.

    \vspace{10pt}
    \begin{theorem}[Teorema di Rouché-Capelli] 
        Sia $S$ un sistema lineare e sia $(A | B)$ la sua matrice aumentata coniugata. $S$ ha soluzione se e solo se 
        $\textnormal{rank}(A | B)$ = $\textnormal{rank}(A)$ e la dimensione dello spazio delle soluzioni di $S$ vale $n - p$,
        con $n$ il numero di incognite del sistema e $p = \textnormal{rank}(A) = \textnormal{rank}((A | B))$
    \end{theorem}

    \begin{proof}[Dimostrazione] 
        Il sistema $S$ può essere scritto in forma: $A\mathbf{x} = \mathbf{b}$. Si riscriva la relazione come: $L_A(\mathbf{x}) = \mathbf{b}$ dove $L_A: \mathbb{R}^n \rightarrow \mathbb{R}^m$ è un'applicazione lineare.
        $S$ ammette soluzione se e solo se $\mathbf{b} \in \textnormal{Im}(L_A)$, ovvero $\exists \mathbf{x} \in \mathbb{R}^n \; | \; L_A(\mathbf{x}) = \mathbf{b}$. E' possibile osservare come l'immagine di $L_A(\mathbf{x})$
        sia data da una combinazione lineare dei vettori colonna di $A$:
        \[
            A\mathbf{x} = 
            \begin{bmatrix}
                a_{11}x_{1} + a_{12}x_{2} + \cdots + a_{1n}x_{n} \\
                a_{21}x_{1} + a_{22}x_{2} + \cdots + a_{2n}x_{n} \\
                \vdots \\
                a_{m1}x_{1} + a_{m2}x_{2} + \cdots + a_{mn}x_{n} \\
            \end{bmatrix}
            = 
            \begin{bmatrix}
                a_{11} \\
                a_{21} \\
                \vdots \\
                a_{m1} \\
            \end{bmatrix}
            x_{1}
            + 
            \begin{bmatrix}
                a_{12} \\
                a_{22} \\
                \vdots \\
                a_{m2} \\
            \end{bmatrix}
            x_{2}
            +
            \cdots
            +
            \begin{bmatrix}
                a_{1n} \\
                a_{2n} \\
                \vdots \\
                a_{mn} \\
            \end{bmatrix}
            x_{n}
        \]    

        \vspace{10pt}
        Il che significa che $\mathbf{b}$ è contenuto nell'immagine di $L_A$ se e solo se $\mathbf{b}$ è nello span delle colonne di $A$.
        Si confronti il rango della matrice aumentata del sistema con il rango della matrice dei coefficienti. Dato che $(A | B)$ è un'estensione di $A$, vi sono due casi: 
        \begin{itemize}
            \item $\textnormal{rank}{(A | B)} = \textnormal{rank}(A)$ 
            \item $\textnormal{rank}{(A | B)} > \textnormal{rank}(A)$ 
        \end{itemize}
        Se $\textnormal{rank}{(A | B)} > \textnormal{rank}(A)$, allora $\mathbf{b}$ è linearmente indipendente ripetto a tutti i vettori di colonna di $A$, il che significa
        che $\mathbf{b} \notin \textnormal{Im}(L_A)$ e che $S$ non ha soluzione. Analogamente, se $\textnormal{rank}{(A | B)} = \textnormal{rank}(A)$, allora $\mathbf{b} \in \textnormal{Im}(L_A)$.

        \vspace{10pt}
        Appurato che esista almeno una soluzione $\mathbf{x_0}$, tutte le restanti soluzioni sono determinabili in maniera seguente:
        \[
            L_A(\mathbf{v}) = \mathbf{0} \implies L_A(\mathbf{x_0} + \mathbf{v}) = L_A(\mathbf{x_0}) + L_A(\mathbf{v}) = \mathbf{b} + \mathbf{0} = \mathbf{b}
        \]
        Di conseguenza, la dimensione dello spazio di soluzione coincide con la dimensione del nucleo di $L_A$, la quale, per il teorema della dimensione, deve rispettare l'uguaglianza:
        \[
            \dim(\textnormal{Im}(L_A)) + \dim(\ker(L_A)) = \dim(\textnormal{Dom}(L_A)) \iff \textnormal{rank}(A) + \textnormal{null}(A) = n
        \] 
        ricavando così la dimensione dello spazio di soluzione come $n - \textnormal{rank}(A)$.  
    \end{proof}




\subsection{Equazioni matriciali}
    Siano $A \in \mathbb{R}^{m, n}, \; X \in \mathbb{R}^{n, k}, \; B \in \mathbb{R}^{m, k}$, allora:
    \[
        AX = B
    \]
    è un'equazione matriciale la cui incognita è $X$.

    \vspace{10pt}
    Per le equazioni matriciali vige il teorema di Rouché-Capelli generalizzato, ovvero: \\
    Una qualsiasi eq. matriciale della forma $AX = B$ ha soluzione se e solo se rank$(A) = $ rank$(A | B)$ e qualora vi fossero delle soluzioni, la cardinalità dello 
    spazio di soluzione è data da $|\mathbb{R}^{m - p, n}|$ dove $p = \textnormal{rank}(A) = \textnormal{rank}(A | B)$.

    \vspace{10pt}
    \begin{corollary}[Corollario di Rouché-Capelli]
        Siano $A, X, B \in \mathbb{R}^{n, n}$ e sia $\textnormal{rank}(B) = n$, allora le seguenti condizioni sono equivalenti:
        \begin{itemize}
            \item $AX = B$ è risolubile
            \item $\textnormal{rank}(A) = n$
            \item La soluzione $X$ è unica 
        \end{itemize}
    \end{corollary}


\vspace{10pt}
\subsection{Matrici inverse, simili, ortogonali e ortonormali}
    Una matrice $A \in \mathbb{R}^{n, n}$ si dice invertibile se e solo se $\exists A^{-1} \in \mathbb{R}^{n, n} \; | \; AA^{-1} = I_n$. Se esiste la matrice inversa, allora essa è sia inversa a sinistra che inversa a destra, secondo quanto segue:
    \[
        C = A^{-1} \; (\textnormal{sx}), \; B = A^{-1} \; (\textnormal{dx}) \implies CA = AB = I \implies C = CI = CAB = IB = B 
    \]

    Due matrici $A, B \in \mathbb{R}^{n, n}$ si dicono simili se:
    \[
        \exists P \in \mathbb{R}^{n, n} \; | \; A = P^{-1}BP  \; \land \; B = P^{-1} 
    \]

    Una matrice $A \in \mathbb{R}^{n, n}$ si dice ortogonale ($A \in O(n)$) se $AA^{T} = I$. $A$ è ortonormale ($A \in SO(n)$) se è ortogonale e $\det(A) = 1$.
    

\end{document}

